% =============================================================================
\chapter{Réalisation — phases A à C}\label{chap:phases-ac}
% =============================================================================

Chaque phase est présentée de la même façon : objectif, réalisation, résultats mesurés (avec la
preuve correspondante), difficultés rencontrées et moyen de retour arrière.

% -----------------------------------------------------------------------------
\section{Phase A — Installation et réseau}\label{sec:phase-a}

\fiche{Mettre en service les machines virtuelles, leurs noms, leurs adresses et la connectivité du laboratoire.}
      {MIR Mahan, Jolan Allen}{19/10/2026}{\fait}

\subsection{Réalisation}

\paragraph{Réseau.} Après sauvegarde de \code{/etc/network/interfaces}, le pont \code{vmbr1}
(\code{lab/proxmox/vmbr1.interfaces}) et une table nftables dédiée
(\code{lab/proxmox/infraready-nat.nft}) ont été ajoutés sur l'hôte : traduction d'adresses
(\emph{masquerade}) du laboratoire vers Internet, et rejet de toute nouvelle connexion du
laboratoire vers \code{192.168.0.0/24}.

\paragraph{Installations automatiques.} Aucune saisie manuelle n'a été nécessaire :
\begin{itemize}
  \item \textbf{LINUX01 et OPS01} : ISO Ubuntu 24.04.5 (empreinte SHA-256 vérifiée) et second lecteur
        \og CIDATA \fg{} portant un fichier \emph{autoinstall} : nom, adresse fixe, LVM, clavier
        français, compte d'administration avec mot de passe haché, SSH par clé uniquement, agent
        invité QEMU, mises à jour de sécurité ;
  \item \textbf{WIN-DC01} : ISO Windows Server 2025 et fichier \code{autounattend.xml} : image Standard
        avec interface graphique, partitionnement UEFI/GPT, nom, adresse fixe, fuseau horaire, WinRM
        et ping autorisés depuis le laboratoire uniquement.
\end{itemize}
Les fichiers de réponse contenant de vrais mots de passe ont été supprimés du stockage Proxmox
après usage ; le dépôt ne contient que des modèles à variables (\code{lab/*.tpl}).

\paragraph{DNS.} Une fois le domaine créé (phase B), les VM Linux interrogent le contrôleur de
domaine (\code{nameservers 192.168.50.10}, domaine de recherche \code{infraready.test}).

\illustration{P-03-win-dc01-installe-server-manager.png}{0.78}{WIN-DC01 installé automatiquement : Server Manager (\preuve{P-03})}{fig:p03}

\subsection{Résultats}

\begin{table}[H]
\centering
\caption{Phase A — vérifications et résultats}\label{tab:phase-a}
\begin{tabularx}{\linewidth}{@{}P{2.8cm}LP{2.3cm}@{}}
\toprule
\textbf{Vérification} & \textbf{Résultat observé} & \textbf{Preuve} \\
\midrule
Trois VM & 3 VM \code{running}, snapshot \code{s01_installation_ok} ; Windows installé en environ 10 min, Ubuntu en 6 min & \preuve{P-01}, \preuve{P-03} \\
Adressage & .10 / .20 / .30, passerelle .1, toutes sur \code{vmbr1} & \preuve{P-02}, \preuve{P-02b}, \preuve{P-04c} \\
Connectivité & Ping entre les 3 VM et vers la passerelle ; Internet joignable (HTTP 200) & \preuve{P-02}, \preuve{P-02b} \\
Isolation & \code{192.168.0.254} bloqué depuis le laboratoire mais joignable depuis l'hôte : le blocage vient bien de la règle nftables & \preuve{P-02} \\
DNS & Enregistrements A, PTR et SRV LDAP résolus depuis les VM Linux et le DC & \preuve{P-04}, \preuve{P-04b}, \preuve{P-04c} \\
\bottomrule
\end{tabularx}
\end{table}

\subsection{Difficultés et décisions}

\begin{table}[H]
\centering
\caption{Phase A — difficultés rencontrées}\label{tab:phase-a-diff}
\begin{tabularx}{\linewidth}{@{}P{3.6cm}P{3.6cm}L@{}}
\toprule
\textbf{Problème} & \textbf{Cause} & \textbf{Correction et résultat} \\
\midrule
ISO Ubuntu 26.04 attachée aux VM Linux & Mauvaise ISO à la création & ISO 24.04.5 téléchargée, empreinte vérifiée, remplacée \\
Carte réseau e1000 sur WIN-DC01 & Pilote non garanti sous Windows Server 2025 & Passage en e1000e (même adresse MAC) \\
L'installeur Ubuntu demande une confirmation & Mot-clé \code{autoinstall} absent de la ligne de commande & Réponse envoyée par \code{qm sendkey} \\
Pas d'iptables sur l'hôte & Proxmox VE 9 repose sur nftables & Règles écrites en nftables \\
Enregistrement PTR du DC absent & Zone inverse créée après la promotion & Enregistrement dynamique déclenché par \code{windows_base.yml} (phase D) \\
Mot de passe sudo de l'hôte visible dans l'en-tête d'une preuve & Le script de capture recopiait la commande & Preuve nettoyée, script corrigé, recherche des secrets avant chaque commit \\
\bottomrule
\end{tabularx}
\end{table}

\textbf{Retour arrière :} snapshot \code{s01_installation_ok} des 3 VM, sauvegarde
\code{interfaces.bak-infraready} sur l'hôte.

% -----------------------------------------------------------------------------
\section{Phase B — Administration Windows}\label{sec:phase-b}

\fiche{Rendre l'annuaire et le DNS opérationnels, documenter comptes, groupes, droits et journaux.}
      {MIR Mahan, Jolan Allen}{19/10/2026}{\fait}

\subsection{Réalisation}

Après le snapshot \code{s02_avant_adds}, toutes les actions ont été lancées depuis OPS01 par WinRM :
\begin{itemize}
  \item \textbf{Forêt AD} (\code{windows_adds.yml}) : forêt \code{infraready.test}, NetBIOS
        \code{INFRAREADY}, niveau fonctionnel Windows 2025, DNS intégré à AD avec redirecteurs, NTP
        externe (indispensable à Kerberos), mot de passe DSRM lu dans le vault chiffré ;
  \item \textbf{Objets de l'annuaire} (\code{scripts/ad-objets.ps1}, rejouable) : unités
        d'organisation, deux groupes globaux de sécurité et quatre comptes, avec changement de mot de
        passe obligatoire à la première connexion.
\end{itemize}

\begin{table}[H]
\centering
\caption{Groupes, membres et droits attribués (principe du moindre privilège)}\label{tab:ad}
\begin{tabularx}{\linewidth}{@{}P{3cm}P{3.2cm}L@{}}
\toprule
\textbf{Groupe} & \textbf{Membres} & \textbf{Droits sur le partage Comptabilité} \\
\midrule
\code{GG_Comptabilite} & l.petit, t.moreau & Modification (partage SMB \og Change \fg{}, NTFS \og Modify \fg{}) \\
\code{GG_Support_IT} & j.martin, s.bernard & Lecture seule, pour dépanner \\
Autres utilisateurs & — & Aucun accès : \texttt{BUILTIN\textbackslash Users} retiré, héritage désactivé, partage masqué \\
\bottomrule
\end{tabularx}
\end{table}

\subsection{Résultats}

\begin{table}[H]
\centering
\caption{Phase B — vérifications et résultats}\label{tab:phase-b}
\begin{tabularx}{\linewidth}{@{}P{2.6cm}LP{2.6cm}@{}}
\toprule
\textbf{Vérification} & \textbf{Résultat observé} & \textbf{Preuve} \\
\midrule
Domaine & \code{infraready.test}, DC \code{WIN-DC01.infraready.test}, catalogue global ; \code{dcdiag} : 10 tests fonctionnels réussis & \preuve{P-10}, \preuve{P-12}, \preuve{P-12b} \\
Comptes et groupes & 2 groupes, 4 comptes, appartenances conformes & \preuve{P-11} \\
DNS interne & A et SRV LDAP résolus, résolution externe par les redirecteurs, zones en mise à jour sécurisée & \preuve{P-12}, \preuve{P-04c} \\
Journaux & Redémarrage de promotion (1074), démarrage d'AD DS (1000), zones DNS chargées, création des comptes et groupes tracée (4720, 4727, 4728) & \preuve{P-13} \\
\bottomrule
\end{tabularx}
\end{table}

\subsection{Difficultés et décisions}

\begin{table}[H]
\centering
\caption{Phase B — difficultés rencontrées}\label{tab:phase-b-diff}
\begin{tabularx}{\linewidth}{@{}P{3.6cm}P{3.8cm}L@{}}
\toprule
\textbf{Problème} & \textbf{Cause} & \textbf{Correction et résultat} \\
\midrule
1\textsuperscript{er} lancement de \code{windows_adds.yml} en échec & Paramètre \code{reboot_timeout} absent de \code{microsoft.ad} 1.4.1 & Paramètre retiré ; forêt créée au 2\textsuperscript{e} lancement (\preuve{P-10a} → \preuve{P-10}) \\
DC synchronisé sur sa propre horloge & Configuration par défaut & NTP externe : strate 2 (\preuve{P-12b}) \\
Événements horodatés dans le futur & Hypothèse : saut d'horloge pendant la promotion, avant NTP & Heure vérifiée identique à OPS01 ; plus de dérive \\
« Le DC utilise son propre DNS » signale un changement à chaque passage & \code{win_dns_client} retire \code{::1} que Windows remet sur un DC & Liste \code{127.0.0.1} et \code{::1} : 2\textsuperscript{e} passage \code{changed=0} (\preuve{P-15}) \\
\bottomrule
\end{tabularx}
\end{table}

\textbf{Retour arrière :} snapshots \code{s02_avant_adds} et \code{s03_avant_ansible} de WIN-DC01.

% -----------------------------------------------------------------------------
\section{Phase C — Administration Linux}\label{sec:phase-c}

\fiche{Gérer identités, permissions, service, stockage, mises à jour et journaux sur Ubuntu.}
      {MIR Mahan, Jolan Allen}{19/10/2026}{\fait}

\subsection{Réalisation}

Réalisée d'abord à la main sur LINUX01 (commandes conservées dans les preuves), puis reprise par
Ansible en phase D.

\begin{table}[H]
\centering
\caption{Comptes, groupes et permissions sur LINUX01}\label{tab:linux-droits}
\begin{tabularx}{\linewidth}{@{}P{3.4cm}P{3.6cm}L@{}}
\toprule
\textbf{Élément} & \textbf{Configuration} & \textbf{Effet} \\
\midrule
alice, bruno & groupe \code{webadmins} & Mettent à jour le site web \\
chloe & groupes \code{auditeurs}, \code{systemd-journal} & Lit les journaux et rapports, sans droit d'écriture \\
ansible & sudo sans mot de passe, clé SSH & Compte de service de l'automatisation \\
\code{/var/www/infraready} & \code{root:webadmins}, mode 2775 (setgid) & Les webadmins écrivent, les fichiers héritent du groupe \\
\code{/data/rapports} & \code{root:auditeurs}, mode 2750 & Les auditeurs lisent, les autres n'ont aucun accès \\
\bottomrule
\end{tabularx}
\end{table}

\begin{itemize}
  \item \textbf{Service} : nginx 1.24 avec un site dédié et des journaux séparés ; \code{nginx -t} avant
        chaque rechargement.
  \item \textbf{Stockage} : disque de 10 Go ajouté à chaud et identifié par son \textbf{numéro de série}
        (\code{INFRA-DATA01}), et non par un nom de périphérique supposé. Le formatage n'a lieu que si cinq
        contrôles sont conformes (taille, aucune signature, aucune partition, non monté, différent du
        disque système). Système de fichiers ext4 monté sur \code{/data} par UUID avec \code{nofail}.
  \item \textbf{Mise à jour contrôlée} : snapshot, pré-contrôles, simulation \code{apt-get -s upgrade},
        mise à jour sans suppression de paquet, contrôle de l'historique apt et du besoin de redémarrage.
\end{itemize}

\subsection{Résultats}

\begin{table}[H]
\centering
\caption{Phase C — vérifications et résultats}\label{tab:phase-c}
\begin{tabularx}{\linewidth}{@{}P{2.6cm}LP{2.3cm}@{}}
\toprule
\textbf{Vérification} & \textbf{Résultat observé} & \textbf{Preuve} \\
\midrule
Permissions & alice écrit dans le site, chloe est refusée ; chloe lit le journal et \code{/data/rapports}, alice est refusée & \preuve{P-20}, \preuve{P-22} \\
Service web & Actif et activé, port 80 en écoute, HTTP 200 depuis OPS01 ; arrêt → HTTP 000, redémarrage → nouveau processus maître & \preuve{P-21} \\
Stockage & Contrôles conformes, \code{/dev/sdb} ext4 monté sur \code{/data} par UUID, remontage par fstab réussi & \preuve{P-22} \\
Mises à jour & 20 paquets mis à jour, code retour 0, pas de redémarrage requis & \preuve{P-23} \\
Journaux & Tentative SSH \og pirate \fg{} refusée, requête 404 tracée, ajout à chaud du disque tracé, aucune erreur & \preuve{P-24} \\
\bottomrule
\end{tabularx}
\end{table}

\attention{Difficulté : chemin de disque deviné}{%
Les premiers contrôles avant formatage visaient un chemin \code{/dev/disk/by-id/...} supposé, qui
n'existait pas. Aucun formatage n'avait lieu dans ce premier script. Correction : lecture du nom
réel et arrêt automatique si un contrôle n'est pas conforme. Le formatage a ainsi visé uniquement
le bon disque.}

\textbf{Retour arrière :} snapshot \code{s02_avant_maj_c} (disque de test compris), sauvegarde de
la configuration réseau dans \code{/root/}.
